% Options for packages loaded elsewhere
\PassOptionsToPackage{unicode}{hyperref}
\PassOptionsToPackage{hyphens}{url}
%
\documentclass[
  brazilian,
]{article}
\usepackage{lmodern}
\usepackage{amssymb,amsmath}
\usepackage{ifxetex,ifluatex}
\ifnum 0\ifxetex 1\fi\ifluatex 1\fi=0 % if pdftex
  \usepackage[T1]{fontenc}
  \usepackage[utf8]{inputenc}
  \usepackage{textcomp} % provide euro and other symbols
\else % if luatex or xetex
  \usepackage{unicode-math}
  \defaultfontfeatures{Scale=MatchLowercase}
  \defaultfontfeatures[\rmfamily]{Ligatures=TeX,Scale=1}
\fi
% Use upquote if available, for straight quotes in verbatim environments
\IfFileExists{upquote.sty}{\usepackage{upquote}}{}
\IfFileExists{microtype.sty}{% use microtype if available
  \usepackage[]{microtype}
  \UseMicrotypeSet[protrusion]{basicmath} % disable protrusion for tt fonts
}{}
\makeatletter
\@ifundefined{KOMAClassName}{% if non-KOMA class
  \IfFileExists{parskip.sty}{%
    \usepackage{parskip}
  }{% else
    \setlength{\parindent}{0pt}
    \setlength{\parskip}{6pt plus 2pt minus 1pt}}
}{% if KOMA class
  \KOMAoptions{parskip=half}}
\makeatother
\usepackage{xcolor}
\IfFileExists{xurl.sty}{\usepackage{xurl}}{} % add URL line breaks if available
\IfFileExists{bookmark.sty}{\usepackage{bookmark}}{\usepackage{hyperref}}
\hypersetup{
  pdftitle={Parte III: Processamento de Linguagem Natural para Linguistas},
  pdfauthor={CiberExt 26-29 · FEELT38103 · Universidade Federal de Uberlândia},
  pdflang={pt-BR},
  hidelinks,
  pdfcreator={LaTeX via pandoc}}
\urlstyle{same} % disable monospaced font for URLs
\setlength{\emergencystretch}{3em} % prevent overfull lines
\providecommand{\tightlist}{%
  \setlength{\itemsep}{0pt}\setlength{\parskip}{0pt}}
\setcounter{secnumdepth}{-\maxdimen} % remove section numbering
\ifxetex
  % Load polyglossia as late as possible: uses bidi with RTL langages (e.g. Hebrew, Arabic)
  \usepackage{polyglossia}
  \setmainlanguage[]{brazilian}
\else
  \usepackage[shorthands=off,main=brazilian]{babel}
\fi

\title{Parte III: Processamento de Linguagem Natural para Linguistas}
\usepackage{etoolbox}
\makeatletter
\providecommand{\subtitle}[1]{% add subtitle to \maketitle
  \apptocmd{\@title}{\par {\large #1 \par}}{}{}
}
\makeatother
\subtitle{Unidade 4 - Exercícios}
\author{CiberExt 26-29 · FEELT38103 · Universidade Federal de
Uberlândia}
\date{Agosto de 2026}

\begin{document}
\maketitle

\hypertarget{exercuxedcios-introduuxe7uxe3o-aos-word-embeddings}{%
\section{\texorpdfstring{Exercícios --- Introdução aos \emph{word
embeddings}}{Exercícios --- Introdução aos word embeddings}}\label{exercuxedcios-introduuxe7uxe3o-aos-word-embeddings}}

Os exercícios abaixo pressupõem que você tenha executado o código da
unidade e que as seguintes variáveis estejam disponíveis: \texttt{nlp},
\texttt{docs}, \texttt{df}, \texttt{lemma\_df}, \texttt{unique\_lemmas},
\texttt{lemma\_vectors}, \texttt{pairs}, \texttt{targets},
\texttt{context} e \texttt{embedding\_model}.

\hypertarget{parte-a-conceitos}{%
\subsection{Parte A --- Conceitos}\label{parte-a-conceitos}}

\textbf{Exercício 1.} Enuncie a hipótese distribucional em uma frase e
explique de que modo a citação de Firth (``\emph{You shall know a word
by the company it keeps}'') a antecipa.

\textbf{Exercício 2.} Explique a diferença entre os eixos
\textbf{sintagmático} e \textbf{paradigmático} de Saussure. Na matriz
criada na seção 2.1, o que representam as linhas e o que representam as
colunas em relação a esses eixos?

\textbf{Exercício 3.} Por que a representação obtida na seção 2.1 é
chamada de ``saco de palavras'' (\emph{bag of words})? Cite duas
limitações dessa representação.

\textbf{Exercício 4.} Qual é a diferença entre um vetor \textbf{esparso}
e um vetor \textbf{denso}? Por que os \emph{word embeddings} aprendidos
pela camada oculta são densos?

\textbf{Exercício 5.} Explique por que a predição da palavra vizinha é
descrita como uma \textbf{tarefa substituta} (\emph{proxy}) e não como o
objetivo real do treinamento.

\textbf{Exercício 6.} A camada oculta da rede tem apenas dois neurônios.
Explique por que ela é descrita como um ``gargalo'' e qual o papel desse
gargalo no aprendizado das representações.

\hypertarget{parte-b-programauxe7uxe3o}{%
\subsection{Parte B --- Programação}\label{parte-b-programauxe7uxe3o}}

\textbf{Exercício 7.} Calcule a similaridade de cossenos entre os
\emph{Docs} de índices 3 e 4 usando a matriz \texttt{df}. O valor obtido
é coerente com o conteúdo das duas orações? Justifique.

\textbf{Exercício 8.} Acrescente uma sexta oração à lista
\texttt{examples} --- por exemplo,
\texttt{"Stockholm\ is\ the\ capital\ of\ Sweden"} --- e refaça o
pipeline até a matriz de similaridade de cossenos. Quantas colunas passa
a ter o \emph{DataFrame} \texttt{df}? Qual \emph{Doc} se torna mais
similar ao novo?

\textbf{Exercício 9.} Usando a matriz de coocorrência
\texttt{lemma\_df}, calcule a similaridade de cossenos entre os vetores
de \texttt{\textquotesingle{}capital\textquotesingle{}} e
\texttt{\textquotesingle{}the\textquotesingle{}}, e entre
\texttt{\textquotesingle{}Helsinki\textquotesingle{}} e
\texttt{\textquotesingle{}ferry\textquotesingle{}}. Comente os
resultados à luz da hipótese distribucional.

\textbf{Exercício 10.} Modifique o laço que preenche \texttt{lemma\_df}
para considerar uma janela de \textbf{três} palavras de cada lado, em
vez de duas. Recalcule a similaridade entre
\texttt{\textquotesingle{}Helsinki\textquotesingle{}} e
\texttt{\textquotesingle{}Tallinn\textquotesingle{}} e compare com o
valor obtido na unidade.

\textbf{Exercício 11.} Escreva uma função
\texttt{one\_hot(lemma,\ vocabulary)} que receba um lema e uma lista de
lemas e devolva o vetor one-hot correspondente, sem usar o dicionário
\texttt{lemma\_vectors}. A função deve levantar um \texttt{ValueError}
se o lema não estiver no vocabulário.

\textbf{Exercício 12.} Treine novamente a rede neural, desta vez com
\textbf{quatro} neurônios na camada oculta, mantendo os demais
parâmetros. Verifique o formato da matriz de pesos da camada oculta e
explique por que já não é possível plotar diretamente o espaço de
embeddings em duas dimensões.

\textbf{Exercício 13.} A partir da matriz
\texttt{hidden\_layer\_weights} obtida na unidade, calcule a
similaridade de cossenos entre os embeddings de
\texttt{\textquotesingle{}Helsinki\textquotesingle{}} e
\texttt{\textquotesingle{}Tallinn\textquotesingle{}}. Compare esse valor
com o obtido na seção 2.2 a partir da matriz de coocorrência.

\hypertarget{parte-c-reflexuxe3o}{%
\subsection{Parte C --- Reflexão}\label{parte-c-reflexuxe3o}}

\textbf{Exercício 14.} O exemplo da unidade usa apenas cinco orações.
Liste três consequências concretas dessa escassez de dados para a
qualidade dos embeddings aprendidos.

\textbf{Exercício 15.} A unidade menciona que abordagens contemporâneas
tentam distinguir formas homonímicas como ``banco'' (instituição) e
``banco'' (assento). Explique por que o modelo construído nesta unidade
é \textbf{incapaz} de fazer essa distinção e o que seria necessário para
superá-la.

\end{document}
